%% ============================================================
%%  生成AI利用申告書 — 提出課題1 (HCK02_01.ino)
%%  ハッカソン1 / 塩澤 拓海 / 2026
%%  コンパイル: docker run --rm -v "$(pwd):/workspace" -w /workspace \
%%              ghcr.io/paperist/texlive-ja:debian latexmk main.tex
%% ============================================================
\documentclass[a4paper,11pt,dvipdfmx]{jsarticle}

%% ─── レイアウト ─────────────────────────────
\usepackage[top=25mm,bottom=25mm,left=20mm,right=20mm]{geometry}

%% ─── 色定義（プロジェクト共通パレット）──────
\usepackage{xcolor}
\definecolor{headerblue}{HTML}{1B3A5C}
\definecolor{accentblue}{HTML}{2E86C1}
\definecolor{lightgray}{HTML}{F2F3F4}
\definecolor{codebg}{HTML}{F7F9FB}
\definecolor{warnred}{HTML}{E74C3C}
\definecolor{noteyellow}{HTML}{F39C12}
\definecolor{okgreen}{HTML}{27AE60}

%% ─── 見出し ─────────────────────────────────
\usepackage{titlesec}
\usepackage{needspace}
\titleformat{\section}
  {\needspace{5\baselineskip}\Large\bfseries\color{headerblue}}
  {\thesection}{1em}{}
  [\vspace{-0.5em}\textcolor{accentblue}{\rule{\textwidth}{1.2pt}}]
\titleformat{\subsection}
  {\needspace{4\baselineskip}\large\bfseries\color{headerblue}}
  {\thesubsection}{1em}{}

%% ─── 改ページ制御 ───────────────────────────
\widowpenalty=10000
\clubpenalty=10000
\predisplaypenalty=10000
\interfootnotelinepenalty=10000
\brokenpenalty=5000

%% ─── tcolorbox ─────────────────────────────
\usepackage[most]{tcolorbox}
\tcbuselibrary{skins,breakable}
\tcbset{before={\par\medskip\needspace{2\baselineskip}\noindent}}

\newtcolorbox{itembox}[1]{%
  enhanced, breakable,
  colback=codebg, colframe=accentblue,
  boxrule=0.5pt, arc=2pt,
  left=12pt, right=12pt, top=10pt, bottom=10pt,
  fonttitle=\bfseries, coltitle=white, colbacktitle=accentblue,
  attach boxed title to top left={xshift=10pt, yshift=-3mm},
  boxed title style={sharp corners, boxrule=0pt, arc=2pt,
                     left=6pt, right=6pt, top=2pt, bottom=2pt},
  title={#1}
}

%% ─── ヘッダー／フッター ─────────────────────
\usepackage{fancyhdr}
\pagestyle{fancy}
\fancyhf{}
\fancyhead[L]{\small\color{headerblue}生成AI利用申告書}
\fancyhead[R]{\small\color{headerblue}\assignname}
\fancyfoot[C]{\small\thepage}
\renewcommand{\headrulewidth}{0.4pt}

%% ─── リスト・表 ─────────────────────────────
\usepackage{booktabs}
\usepackage{tabularx}
\usepackage{longtable}
\usepackage{enumitem}
\setlist{nosep, leftmargin=1.5em}

%% ─── チェックボックス ───────────────────────
\usepackage{amssymb}
\newcommand{\cboxon}{$\boxtimes$\ }
\newcommand{\cboxoff}{$\square$\ }

%% ─── コードブロック（listings + plistings）──
\usepackage{listings}
\usepackage{plistings}

\lstdefinestyle{promptstyle}{
  basicstyle=\ttfamily\small,
  backgroundcolor=\color{codebg},
  frame=single, framerule=0.4pt, rulecolor=\color{accentblue!60},
  breaklines=true, breakatwhitespace=false,
  xleftmargin=8pt, xrightmargin=8pt,
  aboveskip=6pt, belowskip=6pt,
  columns=flexible,
}
\lstdefinestyle{cppstyle}{
  language=C++,
  basicstyle=\ttfamily\footnotesize,
  keywordstyle=\color{accentblue}\bfseries,
  commentstyle=\color{gray}\itshape,
  stringstyle=\color{okgreen},
  numberstyle=\tiny\color{gray},
  numbers=left, numbersep=8pt,
  frame=single, framerule=0.5pt, rulecolor=\color{accentblue},
  backgroundcolor=\color{codebg},
  breaklines=true, breakatwhitespace=false,
  tabsize=2, showstringspaces=false,
  xleftmargin=15pt, framexleftmargin=15pt,
  morekeywords={uint8_t,uint16_t,uint32_t,int16_t,int32_t,setup,loop,analogRead,Serial},
}

%% ─── 課題メタ情報 ───────────────────────────
\newcommand{\assignname}{提出課題1 (HCK02\_01.ino)}

%% ─── hyperref ───────────────────────────────
\usepackage[dvipdfmx]{hyperref}
\usepackage{pxjahyper}
\hypersetup{
  colorlinks=true,
  linkcolor=accentblue, urlcolor=accentblue, citecolor=accentblue,
  pdftitle={生成AI利用申告書 提出課題1},
  pdfauthor={Shiozawa Takumi}
}

%% ─── TODO 強調マクロ ────────────────────────
\newcommand{\todo}[1]{\textcolor{warnred}{\textbf{[TODO: #1]}}}

%% ============================================================
\begin{document}

%% ─── タイトル ───────────────────────────────
\begin{center}
  {\LARGE\bfseries\color{headerblue}生成AI利用申告書}\\[2pt]
  {\small\color{accentblue}\assignname}
\end{center}
\vspace{2pt}
\textcolor{accentblue}{\rule{\textwidth}{1.2pt}}
\vspace{6pt}

%% ─── 基本情報 ───────────────────────────────
\noindent
\begin{tabularx}{\textwidth}{@{}l X@{}}
  \textbf{氏\ 名}               & 塩澤 拓海 \\
  \textbf{学籍番号}             & \todo{学籍番号を記入} \\
  \textbf{プログラムファイル名} & \texttt{HCK02\_01.ino} \\
  \textbf{使用したAIツール}     & Claude Code (Opus 4.7) \\
\end{tabularx}

\vspace{10pt}

%% ============================================================
%%  1. 何に使ったか
%% ============================================================
\begin{itembox}{1.\ 何に使ったか}
例題2 (\texttt{work3.ino}) のシリアル出力プログラムについて、
以下の3点で生成AIを利用した。
\begin{itemize}
  \item シリアルプロッタで音声波形がきれいに描画されない原因の診断
  \item 修正方針（ボーレート・DCオフセット除去・サンプリング間隔・出力ラベル）の洗い出し
  \item 修正後コードのコメント推敲
\end{itemize}
\end{itembox}

%% ============================================================
%%  2. AIへの指示（プロンプト）
%% ============================================================
\begin{itembox}{2.\ AIへの指示（プロンプト）}
実際にAIへ送った指示のうち、主要なものを以下に示す。

\begin{lstlisting}[style=promptstyle]
例題2 (work3.ino) のコードを読んで、Arduinoのシリアルプロッタに
音声波形がきれいに表示されない原因を診断してほしい。

以下の観点で問題点を洗い出して、修正方針を箇条書きで出して:
- ボーレート設定
- DCオフセット除去の仕方
- サンプリング間隔
- シリアルプロッタに出力するデータの種類とスケール
- 変数名やコメントで直すべきところ
\end{lstlisting}

\begin{lstlisting}[style=promptstyle]
一緒に HCK02_01.ino として修正版を作りたい。
例題2 からの変更点をコメントで明示して、
あとで自分の言葉で説明できる形にしてほしい。
\end{lstlisting}
\end{itembox}

%% ============================================================
%%  3. AIの出力の使い方
%% ============================================================
\begin{itembox}{3.\ AIの出力の使い方}
\cboxoff そのまま使った\quad
\cboxon 一部修正して使った\quad
\cboxoff 参考にしただけ

\vspace{6pt}
\textbf{選んだ理由：}
AIが挙げた原因と修正方針のうち、以下は妥当と判断して \texttt{HCK02\_01.ino} に採用した：
\begin{itemize}
  \item 実測値 (0--1023) と振幅 (V単位) を同時に出力しているため、
        プロッタのY軸が整数側に引っ張られて振幅波形がつぶれる問題
        $\rightarrow$ 実測値の出力を撤去し、振幅系の3トレースだけ出す
  \item DCオフセットを固定値 (1.65V) で引いているため、
        実際のバイアスが違うと中心が0にならない問題
        $\rightarrow$ EWMA による動的追従 (\texttt{ALPHA\_DC}=0.001) に変更
  \item \texttt{delay} がないためシリアル送信バッファが溢れ、
        波形がガタつく問題
        $\rightarrow$ \texttt{delayMicroseconds(500)} で 2\,kHz サンプリングに固定
\end{itemize}
一方で、AIから以下の追加提案もあったが、\textbf{実機で動作確認していない現段階では採用を保留}している：
\begin{itemize}
  \item ボーレート 921600 $\rightarrow$ 115200 への変更
        （実機で取りこぼしが確認できてから判断する）
  \item 出力ラベルを日本語から英字 (\texttt{signal}/\texttt{ref\_hi}/\texttt{ref\_lo}) へ変更
        （手元環境で文字化けが実際に起きてから判断する）
  \item 参照電圧を \texttt{\#define VREF} で定数化
        （基板を切り替える必要が出たら導入する）
\end{itemize}
EWMA の時定数 (\texttt{ALPHA\_DC}=0.001 / \texttt{ALPHA\_LP}=0.2) は
AIの提案値をそのまま採用せず、実機でマイク回路の特性を見ながら調整する方針にしている。
\end{itembox}

%% ============================================================
%%  4. 正しいと確認した方法
%% ============================================================
\begin{itembox}{4.\ 正しいと確認した方法}
\todo{実機テスト後に具体的な観測内容を記入する}

想定している確認手順（実機テストで実施する）：
\begin{itemize}
  \item Arduino IDE でコンパイルが通ることを確認する
  \item シリアルモニタを 921600\,bps で開き、数値が流れることを確認する
  \item シリアルプロッタに切り替え、「振幅 / 最大振幅 / 最小振幅」の3本が見えることを確認する
  \item 無音時に「振幅」が 0 付近に収束することを観察する（DC 追従の確認）
  \item マイクに声をかけると「振幅」が上下に振動することを観察する
  \item コードを1行ずつ読み、EWMA 式 \(y_n = \alpha x + (1-\alpha) y_{n-1}\) の意味を説明できる状態にする
\end{itemize}
\end{itembox}

%% ============================================================
%%  5. 問題や間違い
%% ============================================================
\begin{itembox}{5.\ 問題や間違い}
\cboxoff あった\quad
\cboxoff なかった

\vspace{6pt}
\todo{実機テスト後に、実際に気づいた問題や AI の誤り・自分が気づいた点を記入する}

\textbf{検証予定ポイント（実機で確認する）：}
\begin{itemize}
  \item EWMA の初期値 \texttt{dc = 512.0f} が、
        起動直後に実際の DC バイアスに収束するまで何秒かかるか
  \item \texttt{ALPHA\_DC=0.001} で DC 追従が遅すぎて、
        長めの声など低周波成分まで消してしまわないか
  \item ボーレート 921600 + \texttt{delayMicroseconds(500)} + 3 トレース出力で
        シリアル送信が間に合うか（取りこぼしが起きるようなら 115200 へ下げることを検討）
\end{itemize}
\end{itembox}

%% ============================================================
%%  6. 自分で工夫した点
%% ============================================================
\begin{itembox}{6.\ 自分で工夫した点}
\begin{itemize}
  \item 例題2 が「実測値 (0--1023) と振幅 (V 単位) を同じプロッタに出していた」
        ために波形がつぶれる問題を自分で理解し、振幅系の3トレースのみに絞って出力した
  \item 固定値で DC を引いていたのを EWMA による動的追従に書き換え、
        追従用 (\texttt{ALPHA\_DC}) とノイズ除去用 (\texttt{ALPHA\_LP}) の
        係数を目的別に分けて命名した
  \item \texttt{delayMicroseconds(500)} でサンプリング間隔を 2\,kHz に固定し、
        送信タイミングを安定させた
  \item 「最大振幅」「最小振幅」を定数として常に出力することで、
        プロッタのY軸範囲が振幅レンジ（$\pm$1.65\,V）に張り付くようにした
\end{itemize}
\end{itembox}

%% ============================================================
%%  参考文献
%% ============================================================
\section*{参考文献}
\begin{itemize}
  \item 生成AI利用申告書テンプレート（ハッカソン1 講義資料, \texttt{references/lectures/})
  \item Arduino Reference: \texttt{analogRead}, \texttt{analogReadResolution}, \texttt{Serial}
        \url{https://docs.arduino.cc/language-reference/}
  \item 例題2 オリジナルコード: \texttt{work/shiozawa/work-0422/work3/work3.ino}
\end{itemize}

\end{document}
